% ============================================================
% Chapter 2: Literature Review and Theoretical Background
% ============================================================
\chapter{Literature Review and Theoretical Background}
\label{ch:background}

\chapteroverview{This chapter provides the theoretical and scientific foundation supporting the proposed integrated assistive system. It begins with the medical background necessary to understand the biophysical nature and origin of electrooculography (EOG) signals. It then discusses the fundamental engineering principles related to biomedical signal acquisition, digital signal processing, and physical control systems. Furthermore, the chapter presents a literature review of previous studies in the field of eye-tracking and biomedical assistive technologies. Finally, it concludes by highlighting the advantages of the adopted technical approach and demonstrating the technical feasibility of this integrated design.}


% ============================================================
% Section 2.1: Medical Background
% ============================================================
\section{Medical Background}
\label{sec:medical_background}

\subsection{Anatomy and Physiology of the Human Eye (Corneo-Retinal Dipole)}
\label{subsec:anatomy_eye_dipole}

From an electrical perspective, the human eye functions as a biological dipole or bio-battery. This electrostatic property arises from the resting potential difference maintained between the cornea and the retina. The anterior segment (cornea) carries a positive electrical charge, whereas the posterior segment (retina) carries a negative charge. This charge distribution generates an internal standing potential ranging from $10\,\text{mV}$ to $30\,\text{mV}$ \cite{webster2009medical}. However, when measured at the skin surface, this potential experiences significant attenuation due to the impedance of surrounding biological tissues. As the eyeball rotates, the orientation and magnitude of this electrostatic field shift relative to fixed electrode placement sites on the face. Figure \ref{fig:corneo_retinal_dipole} illustrates this corneo-retinal dipole mechanism and the resulting voltage variation observed at the skin surface during horizontal eye movements.

\begin{figure}[H]
  \centering
  \includegraphics[width=0.75\textwidth]{Graphical-illustration-of-the-corneo-retinal-standing-potential-of-the-eyeball.png}
  \caption{Corneo-retinal dipole potential and signal variations during eye movements \cite{dietrich2017towards}.}
  \label{fig:corneo_retinal_dipole}
\end{figure}

\FloatBarrier
\subsection{Biomechanics of Eye Movements and Blinks}
\label{subsec:biomechanics_eye_movements}

Human ocular dynamics consist primarily of rapid ballistic movements known as saccades and stationary resting intervals referred to as fixations. In addition, eyelid activity (blinking) can be physiologically categorized into involuntary and voluntary blinks. Involuntary blinks occur spontaneously to lubricate the corneal surface, whereas voluntary blinks are intentional, displaying significantly higher signal amplitudes and longer durations \cite{hall2016guyton}. Distinguishing between these two physiological behaviors is essential for the system's control logic. The software algorithm relies on detecting intentional, strong blinks and sustained fixations to generate mode-switching and selection commands while filtering out random involuntary blinks to prevent false system activation.

To distinguish between spontaneous eyelid activity and intentional control commands, the temporal and amplitude characteristics of the biopotential waveforms were analyzed. As illustrated in Figure \ref{fig:eog_blink_thresholding}(A), an involuntary blink produces a narrow, lower-amplitude biphasic spike with a short temporal footprint ($\approx 150\,\text{ms}$) resulting from spontaneous corneal lubrication. Conversely, as depicted in Figure \ref{fig:eog_blink_thresholding}(B), an intentional voluntary blink yields a significantly wider peak with a higher voltage amplitude and an extended temporal duration ($>400\,\text{ms}$) caused by conscious eyelid muscle contraction. These amplitude and temporal width disparities form the physiological basis for real-time command classification.

\begin{figure}[H]
  \centering
  \includegraphics[width=0.95\textwidth]{eog_blink_thresholding.png}
  \caption{Comparison of EOG signal waveforms: (A) Spontaneous involuntary blink displaying a narrow peak and short duration ($\approx 150\,\text{ms}$); (B) Intentional voluntary blink exhibiting a wider peak, higher voltage amplitude, and extended duration ($>400\,\text{ms}$).}
  \label{fig:eog_blink_thresholding}
\end{figure}

\FloatBarrier
\subsection{Characteristics of the Electrooculography (EOG) Signal}
\label{subsec:eog_characteristics_detail}

The electrooculography (EOG) signal possesses specific biophysical characteristics that dictate the design requirements of the signal conditioning hardware. The amplitude of raw surface-measured EOG signals typically ranges from $10\,\mu\text{V}$ to $5\,\text{mV}$, necessitating high-precision differential instrumentation amplifiers. In terms of frequency distribution, the fundamental power spectrum of interest for ocular movement spans from $0.1\,\text{Hz}$ to $30\,\text{Hz}$ \cite{webster2009medical}. Based on these biophysical properties, a hardware low-pass filter with a cut-off frequency of $30\,\text{Hz}$ was implemented to eliminate unwanted high-frequency noise. Furthermore, to isolate the DC offset component and shape the signal into sharp, well-defined pulses (spikes) corresponding to voluntary blinks, the cut-off frequency of the high-pass filter was set to $0.8\,\text{Hz}$.

\FloatBarrier
\subsection{Physiological Principles of Electrode Placement}
\label{subsec:electrode_placement}

Acquiring high-quality EOG signals requires proper anatomical positioning of surface electrodes. To capture the vertical channel (corresponding to vertical eye movements and eyelid blinks), electrodes are attached directly above and below one eye. For the horizontal channel, electrodes are placed at the lateral margins of both eyes, near the outer canthi. This configuration measures differential potential changes as the corneo-retinal dipole rotates toward or away from a given electrode, while a reference electrode positioned on the forehead provides a ground potential to cancel common-mode interference \cite{webster2009medical}. To illustrate the vertical channel electrode configuration described above, Figure \ref{fig:vertical_eog_placement} depicts the anatomical positioning of the supra-orbital and infra-orbital electrodes used in this project.

\begin{figure}[H]
  \centering
  \includegraphics[width=0.6\textwidth]{vertical_eog_electrode_placement.jpeg}
  \caption{Vertical EOG electrode placement \cite{li2021evaluation}.}
  \label{fig:vertical_eog_placement}
\end{figure}

\FloatBarrier
\subsection{Pathophysiology of Target Neuromuscular Disorders}
\label{subsec:target_disorders}

The developed assistive platform is designed for individuals suffering from severe neuromuscular impairment and loss of voluntary motor control. Key target pathologies include Amyotrophic Lateral Sclerosis (ALS), a progressive neurodegenerative disease; Quadriplegia, resulting in total limb paralysis; and Locked-in Syndrome (LIS), where cognitive awareness remains intact despite voluntary muscle loss. Furthermore, the system targets patients with Muscular Dystrophy (MD), a group of diseases causing progressive muscular weakness. In many stages of MD, the extraocular muscle control remains relatively preserved, making EOG an optimal, non-fatiguing physiological pathway for human-machine interaction \cite{hall2016guyton}.

\FloatBarrier
\subsection{Clinical Justification for EOG-Based Human-Machine Interfaces}
\label{subsec:clinical_justification}

The medical justification for selecting EOG signals as a control medium for these patient populations relies on an important anatomical and neurological principle: the cranial nerves responsible for extraocular muscle control—specifically the oculomotor nerve (CN III), trochlear nerve (CN IV), and abducens nerve (CN VI)—are typically spared during the progression of neurodegenerative conditions like ALS and severe paralysis \cite{hall2016guyton}. Because ocular motility remains functionally preserved, voluntary eye movements represent a reliable, non-fatiguing physiological pathway through which severely impaired patients can interact with human-machine interfaces (HMIs) to regain independence in mobility and communication.

% ============================================================
% Section 2.2: Engineering Background
% ============================================================
\section{Engineering Background}
\label{sec:engineering_background}

Designing an EOG-guided 3D-printed prototype wheelchair and virtual speller system requires integrating bio-instrumentation, analog electronics, digital signal processing, embedded control, and distance sensing.

\FloatBarrier
\subsection{Principles of Bio-Signal Acquisition and Ambient Noise}
\label{subsec:biopotential_acquisition}

Acquiring electrooculography (EOG) signals from the skin surface requires overcoming several physical and electrical challenges. First, high skin-electrode interface impedance significantly attenuates the microvolt-level biological signal. Second, the bio-potential is highly susceptible to ambient electromagnetic interference, primarily $50\,\text{Hz}$ power line noise, which the human body picks up like an antenna and presents as a common-mode voltage across the measurement electrodes. Additionally, motion artifacts resulting from lead wire movement or adjacent electromyographic (EMG) muscle contractions contaminate the raw signal \cite{webster2009medical}. Overcoming these acquisition challenges requires designing a front-end amplifier with high input impedance ($>10^9\,\Omega$) and high common-mode rejection ratio ($\text{CMRR} > 100\,\text{dB}$) prior to any digital processing stage. To visualize the severity of these noise contamination sources, Figure \ref{fig:eog_noise_comparison} presents a conceptual comparison between a raw EOG signal corrupted by 50 Hz interference and baseline drift versus the ideal clean blink waveform that the signal conditioning chain aims to recover.

\begin{figure}[H]
  \centering
  \includegraphics[width=0.9\textwidth]{eog_noise_comparison.png}
  \caption{Conceptual comparison between an ideal clean EOG blink transient and a raw bio-signal contaminated with 50 Hz power-line interference and baseline drift.}
  \label{fig:eog_noise_comparison}
\end{figure}

\FloatBarrier
\subsection{Instrumentation Amplifiers and the Significance of CMRR}
\label{subsec:instrumentation_amp_cmrr}

To isolate weak bio-potentials from severe ambient interference, an instrumentation amplifier (InAmp) is utilized in the primary acquisition stage rather than standard operational amplifiers (Op-Amps). An instrumentation amplifier operates on differential amplification principles; it amplifies the minute differential voltage existing between two measurement electrodes (representing ocular displacement) while virtually rejecting the common-mode potential present identically at both electrodes (such as $50\,\text{Hz}$ power line noise) \cite{webster2009medical}.

The efficiency of an amplifier in suppressing common-mode noise is mathematically defined by the Common-Mode Rejection Ratio ($\text{CMRR}$), expressed in decibels ($\text{dB}$) as:
\begin{equation}
\text{CMRR (dB)} = 20 \log_{10} \left( \frac{A_d}{A_{cm}} \right)
\label{eq:cmrr_formula}
\end{equation}
where $A_d$ represents the differential voltage gain and $A_{cm}$ represents the common-mode voltage gain.

In precision biomedical applications such as EOG acquisition, maintaining a high $\text{CMRR}$ is essential to ensure signal integrity. This engineering requirement justifies selecting the AD620 integrated circuit (IC) for the primary pre-amplification stage in this project's hardware design. The AD620 provides an exceptionally high input impedance and a minimum $\text{CMRR}$ of $100\,\text{dB}$ \cite{ad620datasheet}.

In addition to noise rejection capabilities, the AD620 instrumentation amplifier allows precise control over the differential voltage gain ($G$) using a single external resistor ($R_G$), governed by the manufacturer's characteristic transfer equation \cite{ad620datasheet}:
\begin{equation}
G = 1 + \frac{49.4\,\text{k}\Omega}{R_G}
\label{eq:ad620_transfer}
\end{equation}
This adjustable feature offers necessary engineering flexibility to achieve the target preliminary amplification for weak blink transients prior to passing the signal through subsequent active filtering stages. To illustrate this integrated architecture, Figure \ref{fig:ad620_functional_schematic} depicts the internal functional schematic of the AD620 instrumentation amplifier and the external connection of the gain-setting resistor ($R_G$) across Pins 1 and 8 \cite{ad620datasheet}.

\begin{figure}[H]
  \centering
  \includegraphics[width=0.65\textwidth]{ad620.jpeg}
  \caption{Functional architecture and $R_G$ gain-setting connection of the AD620 InAmp \cite{ad620datasheet}.}
  \label{fig:ad620_functional_schematic}
\end{figure}

\FloatBarrier
\subsection{Active Analog Filters (Active Sallen-Key and Notch Filters)}
\label{subsec:active_analog_filters}

Active analog filters are critical for conditioning acquired bio-potentials prior to analog-to-digital conversion. Active topologies incorporating operational amplifiers are preferred in low-frequency biomedical applications because they eliminate the need for bulky inductors and offer low output impedance, which prevents inter-stage loading effects \cite{webster2009medical}.

The Sallen-Key architecture is one of the most widely implemented topologies for designing second-order low-pass and high-pass filters. Assuming equal component values ($R_1 = R_2 = R$ and $C_1 = C_2 = C$), the cut-off frequency ($f_c$) is governed by:
\begin{equation}
f_c = \frac{1}{2 \pi R C}
\label{eq:sallen_key_fc_bg}
\end{equation}

In biomedical signal processing, a high-pass filter (HPF) removes the DC offset caused by half-cell potentials at the skin-electrode interface, whereas a low-pass filter (LPF) restricts the signal bandwidth to satisfy sampling requirements and prevent aliasing. Furthermore, to eliminate $50\,\text{Hz}$ or $60\,\text{Hz}$ power-line interference, a notch filter (band-stop filter) is incorporated to sharply attenuate the specific power line frequency without distorting the primary biophysical spectrum. To visualize the spectral selectivity of these conditioning stages, Figure \ref{fig:active_filters_bode} illustrates the frequency response characteristics (Bode plots) of the high-pass ($f_c = 0.8\,\text{Hz}$), low-pass ($f_c = 30\,\text{Hz}$), and 50 Hz notch filters \cite{webster2009medical}.

\begin{figure}[H]
    \centering
    \includegraphics[width=0.65\textwidth]{Active_Filters_EOG_30Hz.png}
    \caption{Ideal frequency response curves (Bode plots) for High-Pass ($f_c = 0.8\,\text{Hz}$), Low-Pass ($f_c = 30\,\text{Hz}$), and 50 Hz Notch active filters in EOG signal conditioning \cite{webster2009medical}.}
    \label{fig:active_filters_bode}
\end{figure}

\FloatBarrier
\subsection{Digital Signal Processing and Pattern Extraction}
\label{subsec:dsp_moving_average}

Once digitized by an Analog-to-Digital Converter (ADC) at a sampling rate adhering to the Nyquist-Shannon sampling theorem, raw bio-signals undergo digital processing to enhance signal quality and extract feature metrics. Processing typically begins with digital smoothing techniques, such as moving average filtering, to suppress high-frequency ripple components and highlight primary signal morphology.

Pattern extraction primarily relies on threshold detection algorithms. Mathematically, the instantaneous discrete signal sample $x(n)$ is compared against a predefined threshold value $Th$. When $x(n) > Th$, the event onset is logged. To differentiate between distinct physiological patterns (such as involuntary vs. intentional blinks), the algorithm measures the elapsed pulse width ($T_{\text{width}}$) during which the signal remains above the threshold. Based on this temporal width and event frequency, the signal is classified to generate real-time control commands for external hardware. To demonstrate this temporal feature extraction logic, Figure \ref{fig:eog_threshold_twidth} depicts the threshold crossing criterion $x(n) > Th$ and the measurement of pulse width ($T_{\text{width}}$) used for pattern classification.

\begin{figure}[H]
    \centering
    \includegraphics[width=0.65\textwidth]{EOG_Threshold_Twidth.png}
    \caption{Conceptual digital thresholding scheme and temporal pulse width ($T_{\text{width}}$) feature extraction for EOG event classification.}
    \label{fig:eog_threshold_twidth}
\end{figure}

\FloatBarrier
\subsection{Motor Control and Pulse-Width Modulation (PWM)}
\label{subsec:dc_motor_pwm}

Guiding robotic and mobility systems requires precise control over both the rotational speed and direction of direct current (DC) motors. Because microcontrollers cannot directly deliver the high currents required by DC motors, motor driver circuits act as high-current electronic switches interfacing the controller to the motor power supply.

To regulate motor speed efficiently, Pulse-Width Modulation (PWM) is employed. PWM rapidly switches the applied voltage between ON and OFF states at a fixed switching frequency. The ratio of the ON duration to the total period, known as the duty cycle ($D$), determines the effective average output voltage ($V_{\text{out}}$) applied to the motor:
\begin{equation}
V_{\text{out}} = D \times V_{\text{in}}
\label{eq:pwm_vout}
\end{equation}
where $V_{\text{in}}$ represents the supply voltage. The duty cycle $D$ is calculated as a percentage:
\begin{equation}
D = \frac{T_{\text{on}}}{T_{\text{total}}} \times 100\%
\label{eq:pwm_duty_cycle}
\end{equation}
where $T_{\text{on}}$ is the active pulse duration and $T_{\text{total}}$ is the total period ($T_{\text{on}} + T_{\text{off}}$). This technique delivers linear speed control while minimizing thermal power dissipation. To illustrate the relationship between pulse active time and effective voltage, Figure \ref{fig:pwm_duty_cycles} compares PWM pulse trains at different duty cycles and their resulting average output voltages.

\begin{figure}[H]
    \centering
    \includegraphics[width=0.65\textwidth]{PWM_Duty_Cycles_Comparison.png}
    \caption{Pulse-Width Modulation (PWM) pulse trains at various duty cycles ($25\%$, $50\%$, and $75\%$) and their corresponding average output voltages.}
    \label{fig:pwm_duty_cycles}
\end{figure}

\FloatBarrier
\subsection{Principle of Ultrasonic Distance Measurement}
\label{subsec:ultrasonic_tof}

To maintain navigation safety and prevent collisions in mobile platforms, ultrasonic sensors are implemented based on the Time-of-Flight (ToF) principle. The transmitter transducer emits ultrasonic acoustic pulses beyond human audible range, while the receiver transducer detects the echo reflected back from an obstacle.

The physical distance ($d$) between the sensor and the detected obstacle is calculated from the elapsed round-trip time ($t$) and the speed of sound in air ($v \approx 343\,\text{m/s}$) using:
\begin{equation}
d = \frac{v \times t}{2}
\label{eq:tof_distance}
\end{equation}
The factor of 2 accounts for the two-way acoustic path (propagation to the obstacle and back). The computed distance feeds directly into safety control logic, executing immediate automated intervention such as emergency stopping to prevent physical collision. To visualize this two-way acoustic path, Figure \ref{fig:ultrasonic_tof_principle} illustrates the Time-of-Flight (ToF) measurement principle between the ultrasonic transducer and an obstacle.

\begin{figure}[H]
    \centering
    \includegraphics[width=0.65\textwidth]{Ultrasonic_ToF_Principle.png}
    \caption{Operating principle of ultrasonic Time-of-Flight (ToF) distance measurement showing acoustic pulse emission, obstacle reflection, and two-way travel path ($2d$).}
    \label{fig:ultrasonic_tof_principle}
\end{figure}

% ============================================================
% Section 2.3: Literature Review
% ============================================================
\section{Literature Review}
\label{sec:literature_review}

This section reviews existing literature on Human-Machine Interfaces (HMIs), smart wheelchairs, and assistive communication platforms. The goal is to analyze technical strengths and limitations, ultimately establishing the research gap addressed by the proposed integrated system.

% ------------------------------------------------------------
\subsection[Scientific Comparison: BCI (EEG) vs. EOG]{Scientific Comparison: Brain-Computer Interfaces (EEG) vs. Electrooculography (EOG)}
\label{subsec:eeg_vs_eog_comparison}

Electroencephalography (EEG) and electrooculography (EOG) represent two prominent bio-potentials utilized for assistive control interfaces \cite{onesensor2014proposals}. Comparing both modalities based on prior literature reveals critical operational differences affecting their real-time feasibility:

\begin{itemize}
    \item \textbf{Signal-to-Noise Ratio (SNR):} EOG signals exhibit a significantly higher SNR compared to EEG signals. Brain potentials are inherently microvolt-level, highly complex, and prone to severe interference from ambient noise and muscular activity \cite{onesensor2014proposals}.
    \item \textbf{Latency:} EEG-based systems typically suffer from relatively high latency due to the intensive signal processing required for pattern classification (e.g., P300 or Motor Imagery) \cite{wolpaw2002bci}. Conversely, EOG signals provide faster response times and well-defined, consistent waveforms that are easily detected in real time.
    \item \textbf{Electrode Count:} EEG interfaces require high-density electrode caps covering the scalp, increasing hardware setup complexity and overall cost. In contrast, EOG systems require only a minimal number of facial surface electrodes placed around the eyes, offering superior user comfort \cite{barea2002eog}.
    \item \textbf{Computational Complexity:} Decoding EEG signals demands advanced digital signal processing algorithms, imposing high computational loads. EOG signals possess distinct feature metrics (e.g., voluntary blink spikes) that can be processed on low-cost embedded microcontrollers \cite{kumar2016eog}.
    \item \textbf{Calibration Requirements:} Both modalities face operational calibration challenges. EOG signals are susceptible to baseline DC drift and ocular fatigue, requiring periodic baseline adjustments. Similarly, EEG systems demand extensive user training due to intra-subject signal non-stationarity.
    \item \textbf{Real-Time Performance and Clinical Comfort:} Voluntary eye movement provides an intuitive control modality that remains functional even in patients with severe paralysis. Although EOG control may sometimes be restricted to discrete command sets, it avoids the continuous cognitive strain associated with operating EEG-based interfaces \cite{hosni2018eeg}.
\end{itemize}

\textbf{Summary:} Literature indicates that EOG is more suitable for real-time control of assistive wheelchairs and communication platforms due to its acquisition simplicity, cost-effectiveness, rapid response, and ability to deliver reliable control with minimal computational and physiological overhead \cite{onesensor2014proposals, hosni2018eeg}.

% ------------------------------------------------------------
\FloatBarrier
\subsection{Biosignal-Based Smart Wheelchairs}
\label{subsec:biosignal_smart_wheelchairs}

Smart wheelchairs have evolved to enhance patient autonomy and reduce caregiver reliance. Prior studies have explored several bio-signal interaction modalities:

\begin{itemize}
    \item \textbf{EEG-Based Systems:} Offered autonomous manual control based on Motor Imagery (MI) \cite{liu2024multimodal}. However, these systems exhibit degraded performance in dynamic environments due to user mental fatigue, command latency, and reliance on expensive instrumentation.
    \item \textbf{EOG-Based Systems:} Utilized ocular movements for direct directional control \cite{barea2002eog}. Prior studies demonstrate that EOG-based systems reduce setup time and deliver precise navigation commands \cite{alzahrani2026design}.
    \item \textbf{Hybrid Systems (EEG/EOG):} Integrated EEG and EOG to overcome single-modality constraints; EOG artifacts (such as double blinks) were repurposed as rapid emergency stop commands, while EEG managed continuous directional steering \cite{hosni2018eeg, hybrid2019wheelchair}.
\end{itemize}

\textbf{Engineering Limitations:} Despite these advancements, many existing smart wheelchairs lack integrated autonomous obstacle avoidance mechanisms, exposing patients to collision risks and necessitating continuous caregiver supervision \cite{smartwheelchair2019obstacle, aiiot2024wheelchair}. Furthermore, high computational complexity hinders deployment on low-cost microcontrollers, while reliance on continuous mental focus imposes a heavy cognitive load on the user.

% ------------------------------------------------------------
\FloatBarrier
\subsection{Virtual Keyboards and Spellers}
\label{subsec:virtual_keyboards_spellers}

To enable communication for severely motor-impaired individuals (e.g., LIS patients), various assistive interface modalities have been developed:

\begin{itemize}
    \item \textbf{P300-Based Systems (EEG):} Rely on matrix spellers with flashing rows and columns, requiring intense visual attention from the user \cite{kabbara2015efficient}. Although widely researched, P300 spellers exhibit low typing speeds, high visual fatigue, and low Information Transfer Rates (ITR).
    \item \textbf{Vision-Based / Eye-Tracking Systems:} Utilize camera-based video-oculography to track gaze direction \cite{eyedriven2026review}. While offering high spatial accuracy, these systems are highly sensitive to ambient lighting variations and require head-pose stability, limiting their operational effectiveness in dynamic outdoor environments.
    \item \textbf{EOG-Based Keyboards:} Operate reliably across variable lighting conditions without optical camera hardware \cite{kumar2016eog, preetha2025realtime}. Character selection is achieved through directional eye movements and confirmed via voluntary blinks, achieving practical communication rates.
\end{itemize}

\textbf{Arabic Support and Usability Limitations:} The vast majority of existing EOG- and EEG-based communication platforms are configured exclusively for Latin or Asian character sets, leaving Arabic language support severely lacking \cite{kabbara2015efficient, alzahrani2026design}. The Arabic language possesses distinct structural and grammatical characteristics that necessitate customized interaction interfaces. Recently, integrating Morse code logic with EOG signals has proven effective for Arabic character generation as a low-cost alternative \cite{alzahrani2026design}. Additionally, many EOG interfaces suffer from the ``Midas Touch'' problem---the inability to differentiate between spontaneous eye movements and intentional control commands---leading to frequent false triggers.

% ------------------------------------------------------------
\FloatBarrier
\subsection{Research Gap}
\label{subsec:research_gap}

Critical analysis of existing literature reveals that current assistive technologies typically provide fragmented solutions: they either focus exclusively on mobility via smart wheelchairs or independently address virtual communication interfaces. Unification of both essential functionalities into a single platform is rarely addressed.

Furthermore, prior studies exhibit common technical limitations:

\begin{enumerate}[label=\textbf{\arabic*.}]
    \item Absence of unified platforms providing dual-mode switching between mobility and communication.
    \item High implementation costs and computational complexity that prevent deployment on low-cost embedded systems.
    \item A severe shortage of Arabic language support in bio-signal-based communication interfaces.
    \item Lack of embedded autonomous navigation safety mechanisms (such as automatic obstacle avoidance) fully integrated with a single EOG control interface.
\end{enumerate}

Based on this review, the primary research gap addressed by this project is formulated as follows:

\begin{tcolorbox}[colback=blue!5!white,colframe=blue!75!black,title=\textbf{Formulated Research Gap}]
\textbf{Research Gap:} The lack of a low-cost, unified assistive platform that effectively integrates motorized wheelchair navigation (featuring automated obstacle avoidance) with Arabic virtual speller communication through a single, non-invasive EOG interaction interface.
\end{tcolorbox}

To provide a structured synthesis of the literature, Table \ref{tab:literature_review_summary} presents a comparative scientific summary comparing prior assistive technologies against the proposed integrated system.

\begin{table}[H]
  \centering
  \caption{Comparative Scientific Summary of Prior Literature vs. Proposed Integrated System}
  \label{tab:literature_review_summary}
  \small
  \begin{tabularx}{\textwidth}{@{} p{2.2cm} X X p{2.0cm} p{2.4cm} @{}}
    \toprule
    \textbf{Modality} & \textbf{Key Advantages} & \textbf{Major Technical Limitations} & \textbf{Target Scope} & \textbf{Arabic \& Safety} \\
    \midrule
    \textbf{EEG / BCI} \cite{liu2024multimodal, wolpaw2002bci} & Direct neural decoding & Low SNR, high latency, complex electrode cap, mental fatigue & Single mode (Mobility/Speller) & Unsupported \\
    \midrule
    \textbf{VOG (Camera)} \cite{eyedriven2026review} & High spatial accuracy & Sensitive to lighting, head posture drift, high compute load & Speller only & Rare / Limited \\
    \midrule
    \textbf{EOG Wheelchair} \cite{barea2002eog, hybrid2019wheelchair} & High SNR, low-cost hardware & Lacks real-time obstacle avoidance, no text speller & Navigation only & No safety override \\
    \midrule
    \textbf{EOG Speller} \cite{kumar2016eog, kabbara2015efficient} & Insensitive to ambient illumination variations & Midas-touch false triggers, non-Arabic script design & Speller only & Lacks native Arabic \\
    \midrule
    \textbf{Proposed System} & High SNR, hybrid MATLAB-Arduino DSP, 3-electrode interface & Relies on voluntary blink temporal width classification & \textbf{Dual-Mode} (Wheelchair + Speller) & \textbf{Full Arabic + Auto Obstacle Escape} \\
    \bottomrule
  \end{tabularx}
\end{table}


% ============================================================
% Section 2.4: Advantages and Feasibility of the Proposed Technique
% ============================================================
\section{Advantages and Feasibility of the Proposed Technique}
\label{sec:advantages_proposed_technique}

The developed system presents an integrated engineering solution that directly addresses the technical and operational gaps identified in prior literature by combining 3D-printed prototype wheelchair navigation and an interactive virtual speller into a single platform. Rather than relying on fragmented, single-purpose assistive devices, the proposed design integrates mobility actuation and text communication into a unified architecture for individuals with severe motor impairments, including quadriplegia and locked-in syndrome (LIS). The key engineering advantages and technical feasibility of this design are outlined below.

\subsection{Functional Integration in a Unified Platform}
\label{subsec:functional_integration}

To overcome the limitations of prior assistive systems that require separate devices for locomotion and communication, the proposed platform integrates both functions into a unified assistive system. Utilizing a single biological acquisition channel, raw electrooculography (EOG) signals are processed and routed to a central MATLAB Graphical User Interface (GUI). The user toggles between ``Wheelchair Mode'' and ``Keyboard Mode'' via a four-consecutive-blink command sequence. This unified architecture eliminates the need for multiple hardware setups, reduces physiological and cognitive workload, and establishes a consolidated control environment for assistive interaction.

\subsection{Blink-Based Control and Reduced Muscular Fatigue}
\label{subsec:blink_based_control}

Unlike vision-based eye-tracking systems that suffer from lighting sensitivity or EEG interfaces that cause cognitive exhaustion, the proposed platform relies exclusively on voluntary blink dynamics (short blinks, long blinks, and distinct blink sequences). To minimize muscular fatigue and reduce the total number of required commands, a state-dependent finite state machine was implemented. The algorithm reinterprets biological control signals depending on the active state of the vehicle or text interface. This state-dependent design provides an immediate, reliable system response with minimal physical exertion, aligning with the limited motor capabilities of LIS patients.

\subsection{Arabic Language Support}
\label{subsec:arabic_language_support}

To provide communication for Arabic speakers, the platform incorporates a virtual keyboard featuring a circular layout divided into six sectors. The user interface automatically scans through the sectors sequentially, allowing the user to select the target sector and subsequent character via voluntary blinks. Beyond basic character input, the speller includes text-editing features, including space insertion, single-character backspace, whole-word deletion, and display clearing. This interaction model bridges the language accessibility gap for Arabic-speaking users.

\subsection{Navigational Safety and Caregiver Override}
\label{subsec:navigational_safety}

Unlike conventional electric wheelchairs that place the full burden of collision avoidance on the user, the 3D-printed wheelchair prototype incorporates an obstacle-avoidance safety routine. Front and rear ultrasonic distance sensors continuously monitor the immediate environment to detect obstacles and trigger emergency auto-stop protocols. Furthermore, to provide operational safety, a mobile application provides a secondary wireless control channel (Caregiver Override) via Bluetooth. This override feature enables caregivers to assume manual control of the wheelchair during emergencies or when user fatigue occurs, providing caregiver intervention alongside motorized navigation.

\subsection{Technical Feasibility and Prototyping}
\label{subsec:technical_clinical_feasibility}

From an engineering perspective, the system demonstrates high feasibility by eliminating the need for expensive commercial bio-signal acquisition equipment. A custom, low-cost Analog Front-End (AFE) circuit was designed specifically to acquire microvolt-level EOG signals, utilizing an AD620 instrumentation amplifier alongside active analog and digital filtering (including a $50\,\text{Hz}$ active notch filter for power line noise rejection). By employing standard, off-the-shelf components---such as Arduino microcontrollers, an HC-05 Bluetooth module for wireless data transmission, and a rechargeable lithium-ion battery system managed by a Battery Management System (BMS)---the design balances signal fidelity against total component cost ($<15$\,USD for the AFE). This modular configuration allows straightforward laboratory replication and further clinical evaluation.


% --- Chapter Conclusion & Transition ---
\vspace{1em}
In conclusion, this chapter synthesized the physiological basis of EOG signals, biopotential conditioning electronics, and relevant academic literature. The established advantages of real-time responsiveness and unified hardware integration form the structural parameters for system design. Building upon these theoretical principles, Chapter 3 details the high-level functional block diagrams, operational mode switching logic, and system flowcharts.
